% Notes of the Montreal seminar, 1970 (cote 129, typed part): archive cote 129, site page 48 (PDF page 49).
% typescript. Draft compared: none.
% AI-transcribed by Claude Opus 5.5 (claude-opus-5-5, adaptive thinking, effort high), Message Batches API, batch mtl129_full (msgbatch_01AGnDKHwy6ATCM7ehUa6263), 2026-11-13; not yet checked by eye.
\scanpage{129/48}
\mspage{\hand{II 20}}

\uline{Preuve} (d'après Oda et Oort).

On a une flèche évidente. Il suffit de considérer le cas où
$G$ est unipotent, et en considérant alors la suite de composition formée
des $\mathrm{Im}\,\mathbb{V}^n_G \subset G$ , on se ramène au cas où $\mathbb{V}_G = 0$ .

Pour un tel groupe $G$ , on a
\[ D^*(G) = \mathrm{Hom}_{k\text{-}Gr}(G,\underset{\to}{\mathbb{W}}) = \mathrm{Hom}_{k\text{-}Gr}(G,\mathbb{G}_a) . \]

Par ailleurs, on a un isomorphisme
\[ \mathrm{hom}_{k\text{-}Gr}(G,\mathbb{G}_a) \cong \mathrm{Ker}\{H^0(G,\underline{O}_G) \xrightarrow{\ \Delta^*-pr_1^*-pr_2^*\ } H^0(G,\underline{O}_G) \otimes_k H^0(G,\underline{O}_G)\} \]
où $\Delta : G \times G \to G$ désigne la multiplication de $G$ . Par le changement de
base $k \to K$ , on a un isomorphisme analogue, d'où l'on déduit le corol-%
laire.

\uline{Corollaire 5.2}

Pour tout p-groupe fini sur $k$ , on a
\[ \mathrm{rang}(G) = p^{\mathrm{long}_W(D^*(G))} . \]

\uline{Preuve}. Il suffit de considérer le cas où $G$ est unipotent, et par le
Corollaire 5.1 on peut supposer $k$ algébriquement clos. Alors $G$ admet
une suite de composition ne faisant intervenir que $\mathbb{Z}/p\mathbb{Z}$ et $\alpha_p$ (2.4).
Or on établit aisément que
\[ D^*(\alpha_p) \sim k = W/pW \quad \text{avec} \quad F = V = 0 \]
\[ D^*(\mathbb{Z}/p\mathbb{Z}) \sim k = W/pW \quad \text{avec} \quad V = 0 \ , \ F(\lambda) = \lambda p \ , \]
et la formule est trivialement vérifiée.

\margin{48}
